\documentclass[11pt,twoside]{article}

\usepackage{amsmath,amssymb,amsthm}
\usepackage{graphicx}
\usepackage{hyperref}
\usepackage[numbers]{natbib}
\usepackage{booktabs}
\usepackage{multirow}
\usepackage[margin=1in]{geometry}
\usepackage{xcolor}
\usepackage{tabularx}
\usepackage{float}

\definecolor{codegreen}{rgb}{0,0.6,0}
\definecolor{codegray}{rgb}{0.5,0.5,0.5}
\definecolor{codepurple}{rgb}{0.58,0,0.82}

\newtheorem{theorem}{Theorem}[section]
\newtheorem{lemma}[theorem]{Lemma}
\newtheorem{corollary}[theorem]{Corollary}
\newtheorem{definition}[theorem]{Definition}

\title{Chapter 4: Dynamic Constraint Adaptation for\\Knowledge Graph Question Answering}
\author{}
\date{}

\begin{document}
\maketitle

% ============================================================================
\section{Introduction}
\label{sec:introduction}
% ============================================================================

Graph-Constrained Reasoning (GCR)~\citep{luo2025-graph-constrained-reasoning}
established that pre-compiling all KG paths into a trie and using
prefix-constrained decoding can force LLMs to produce hallucination-free answers
grounded in a knowledge graph. The core idea is simple: enumerate all valid
reasoning paths from the topic entities to candidate answer entities, build a
trie from these paths, and constrain the LLM's token generation to follow only
paths present in the trie.

GCR reports 91.6\% Hits@1 on WebQSP and 74.6\% on CWQ using a static,
pre-compiled trie containing all paths. While effective, this approach has a
fixed cost: every path is included regardless of relevance. Questions with large
graphs may produce tens of thousands of paths, each contributing to the trie's
memory footprint and offering the LLM more opportunity to select an incorrect
path.

This chapter investigates whether \textbf{dynamic constraint adaptation} ---
modifying the trie structure at pre-decode time or during decoding --- can
reduce path count while maintaining accuracy. We implement and evaluate five
dynamic strategies:

\begin{enumerate}
  \item \textbf{TypeOracle Filtering (v1):} Prune paths before decoding using
        type-checking gates on entity types and relation ranges.
  \item \textbf{Step-wise Decoding (v2):} Expand the trie one hop at a time,
        using the model's partial output to guide the next hop's path set.
  \item \textbf{Post-hoc Validation:} Run baseline decoding, then reject
        predictions that violate TypeOracle constraints.
  \item \textbf{Adaptive Enumeration:} Truncate the BFS path list to a fixed
        budget (30, 100, 500 paths) before trie construction.
  \item \textbf{Ontology-Guided Path Planning (label-reason):} Plan at the type
        level using an ontology reasoner, then enumerate only type-compatible
        entity paths.
\end{enumerate}

We evaluate on two datasets---WebQSP (2-hop) and CWQ (4-hop)---using the
GCR-Meta-Llama-3.1-8B-Instruct model with beam search (\(k=5\)).

% ============================================================================
\section{Experimental Setup}
\label{sec:setup}
% ============================================================================

\subsection{Configuration}

\begin{table}[H]
  \centering
  \caption{Experimental configuration.}
  \label{tab:config}
  \begin{tabular}{ll}
    \toprule
    \textbf{Parameter} & \textbf{Value} \\
    \midrule
    Base Model & rmanluo/GCR-Meta-Llama-3.1-8B-Instruct (8B) \\
    Decoding & Beam search, num\_beams=5 \\
    GPU & NVIDIA RTX 4090 (24 GB VRAM) \\
    GPU Memory (model) & \(\sim\)16--18 GB (FP16) \\
    Dataset & RoG-webqsp (1,628 test), RoG-cwq (1,628 test) \\
    Path Index Length & 4 (both datasets) \\
    Max New Tokens & 256 \\
    Sample Timeout & 120 seconds \\
    \bottomrule
  \end{tabular}
\end{table}

\subsection{Datasets}

\textbf{WebQSP} contains 1,628 test questions over Freebase, with an average
path length of 2 hops and average 2,553 paths per question. Questions are
primarily factual (e.g., ``what is the capital of France?'').

\textbf{CWQ} contains 1,628 test questions over Freebase, with an average path
length of 4 hops and average 1,970 paths per question. Questions involve
compositional reasoning (e.g., ``which country has the most languages
spoken?'').

\subsection{Methods}

\begin{table}[H]
  \centering
  \caption{Methods evaluated.}
  \label{tab:methods}
  \begin{tabular}{lll}
    \toprule
    \textbf{Label} & \textbf{Description} & \textbf{Dynamic Aspect} \\
    \midrule
    Baseline & GCR static trie: all BFS paths pre-compiled & None (static) \\
    Filtered & TypeOracle gates applied before trie construction & Pre-decode pruning \\
    v2 & Step-wise hop-by-hop trie expansion & Mid-decode reconstruction \\
    Validate & Baseline, then reject failing TypeOracle checks & Post-hoc filtering \\
    Adaptive-K & Truncate BFS to \(K\) paths (DFS order) & Pre-decode truncation \\
    Label-reason & Ontology reverse reasoning + constrained DFS & Pre-decode pruning \\
    \bottomrule
  \end{tabular}
\end{table}

\subsection{Bug Fix: Beam Search Configuration}

Prior experiments used a bugged beam search configuration: the HuggingFace
\texttt{generation\_config} was not properly passed to
\texttt{model.generate()}, effectively running greedy decoding despite
\texttt{num\_beams \(>\) 1}. All main results (Sections~\ref{sec:webqsp}
and~\ref{sec:cwq}) use the corrected configuration. The correction provides a
consistent +8--9pp lift (Section~\ref{sec:beam-vs-greedy}).

% ============================================================================
\section{WebQSP Results}
\label{sec:webqsp}
% ============================================================================

\subsection{Main Results}

\begin{table}[H]
  \centering
  \caption{WebQSP results (beam search, \(k=5\)).}
  \label{tab:webqsp-main}
  \begin{tabular}{lrrrr}
    \toprule
    \textbf{Method} & \textbf{N} & \textbf{Hits@1} & \textbf{Avg/q (s)} & \textbf{Avg Paths} \\
    \midrule
    Baseline & 580 & \textbf{89.0\%} & 6.62s & 2,553 \\
    Filtered & 580 & \textbf{88.0\%} & 6.58s & 2,213 \\
    \(\Delta\) & --- & \textbf{-1.0pp} & -0.04s & -340 (13.3\%) \\
    \bottomrule
  \end{tabular}
\end{table}

The baseline 89.0\% is within expected noise of the published GCR result
(91.6\%, reported on a different subset of the same test set). Filtering reduces
path count by 13.3\% while incurring only 1pp accuracy loss. Time savings are
negligible because trie construction dominates over path enumeration.

\subsection{Post-hoc Validation Analysis}

Post-hoc validation applies TypeOracle checks to the model's predicted answer
path rather than pre-pruning the trie. On 1,627 WebQSP questions with greedy
decoding:

\begin{table}[H]
  \centering
  \caption{Post-hoc validation analysis (WebQSP, greedy).}
  \label{tab:validate}
  \begin{tabular}{lr}
    \toprule
    \textbf{Metric} & \textbf{Value} \\
    \midrule
    Wrong predictions caught by TypeOracle & 23.8\% \\
    Correct predictions falsely rejected & 1.1\% \\
    Validation precision & 95.5\% \\
    Validation recall & 23.8\% \\
    \bottomrule
  \end{tabular}
\end{table}

TypeOracle validation catches approximately 1 in 4 incorrect predictions while
misclassifying only 1 in 100 correct predictions. TypeOracle constraints are
precise (high precision) but not comprehensive (low recall)---reflecting the
incompleteness of Freebase's type annotations.

\subsection{Adaptive Enumeration (Greedy Decoding)}

Adaptive methods truncate the BFS path list to a fixed budget. Results use
greedy decoding (pre-beam-fix) and serve as a lower bound:

\begin{table}[H]
  \centering
  \caption{Adaptive enumeration results (WebQSP, greedy).}
  \label{tab:adaptive}
  \begin{tabular}{lrrr}
    \toprule
    \textbf{Budget \(K\)} & \textbf{Hits@1} & \textbf{vs Full} & \textbf{Avg time (s)} \\
    \midrule
    Full (2,553 avg) & 80.6\% & --- & 6.71 \\
    500 & 30.7\% & -49.9pp & 1.51 \\
    100 & 12.8\% & -67.8pp & 0.77 \\
    30 & 7.9\% & -72.7pp & 0.59 \\
    \bottomrule
  \end{tabular}
\end{table}

DFS-order truncation performs catastrophically because BFS orders paths by
discovery order, not by relevance. The first \(K\) paths from DFS are typically
the longest, most roundabout paths involving the first-listed topic entity. This strong
selection bias makes simple path-count truncation ineffective.

% ============================================================================
\section{CWQ Results}
\label{sec:cwq}
% ============================================================================

CWQ is substantially harder than WebQSP due to its 4-hop compositional
questions. The path space is larger, the answer types are more diverse, and the
model must maintain coherent reasoning across more intermediate entities.

\subsection{Baseline and Filtered Results}

\begin{table}[H]
  \centering
  \caption{CWQ results (beam search, \(k=5\)).}
  \label{tab:cwq-main}
  \begin{tabular}{lrrrr}
    \toprule
    \textbf{Method} & \textbf{N} & \textbf{Hits@1} & \textbf{Avg/q (s)} & \textbf{Avg Paths} \\
    \midrule
    Baseline & 500 & \textbf{53.2\%} & 6.10s & 1,970 \\
    Filtered & 100 & \textbf{65.0\%} & 6.22s & 1,764 \\
    \(\Delta\) (100q) & --- & \textbf{-4.0pp} & +0.12s & -206 (10.4\%) \\
    \bottomrule
  \end{tabular}
\end{table}

The baseline accuracy drops sharply from WebQSP (89.0\% \(\to\) 53.2\%). This
35.8pp gap is consistent with the published GCR gap (91.6\% vs 74.6\%).

Filtering on CWQ shows a larger accuracy penalty (-4.0pp vs -1.0pp). Two factors
explain this: (1) deeper graphs have fewer alternative routes to correct
answers, so removing any path is more likely to eliminate the gold path; (2)
CWQ questions often involve abstract or rare entity types that Freebase
annotates sparsely, reducing TypeOracle's precision.

Extrapolating the 4pp delta to 500q scale, filtered accuracy is estimated at
\(\sim\)49\%.

\subsection{Step-wise Decoding (v2) --- Negative Result}

\begin{table}[H]
  \centering
  \caption{v2 results vs baseline.}
  \label{tab:v2}
  \begin{tabular}{lrrrr}
    \toprule
    \textbf{Dataset} & \textbf{N} & \textbf{Baseline} & \textbf{v2} & \(\Delta\) \\
    \midrule
    WebQSP & 1,466 & 80.6\% (greedy) & 54.9\% & -36.7pp \\
    CWQ & 500 & 53.2\% & 32.2\% & -21.0pp \\
    \bottomrule
  \end{tabular}
\end{table}

\textbf{Finding: v2 is catastrophically worse than baseline on both datasets.}
The step-wise approach introduces compounding errors: a wrong entity choice at
hop 1 makes all subsequent hops impossible. On WebQSP's 2-hop graphs, v2 falls
36.7pp below baseline. On CWQ's 4-hop graphs, v2 falls 21.0pp below baseline,
achieving only 32.2\%.

\textbf{Root cause:} The per-hop trie construction suffers from tokenization
misalignment. When the model decodes a partial path, entity boundary detection
is unreliable---the tokenizer splits entity names differently depending on the
surrounding context. This means the prefix-constrained trie cannot reliably
recognize which entity the model has committed to at each step, causing the trie
to reject valid continuations and forcing fallback to unconstrained generation.
This is the same tokenization disalignment problem identified by REL-RAG.

\subsection{Ontology-Guided Path Planning (label-reason)}

The ORT-inspired approach first plans at the type level using an ontology graph,
then constrains entity DFS to follow only type-compatible entities. We construct
a category-level ontology by aggregating 46 Freebase types into 7 broad
categories (Person, Location, Organization, Event, Work, Product, Other) with
24 edges (reduced from 1,134 edges in the cross-product of all Freebase types).

\begin{table}[H]
  \centering
  \caption{Label-reason results (CWQ, 10 questions).}
  \label{tab:label-reason}
  \begin{tabular}{lr}
    \toprule
    \textbf{Metric} & \textbf{Value} \\
    \midrule
    Processable & 9/10 (90\%) \\
    Hits@1 & 3/9 (33.3\%) \\
    Avg time & 3.83s/q \\
    Avg paths constrained & \(\sim\)450 \\
    \bottomrule
  \end{tabular}
\end{table}

The category-level ontology eliminates path explosion, but accuracy is below
baseline (33.3\% vs 53.2\%, evaluated on only 10 questions). The category
abstraction is too coarse: paths valid at the category level may involve entity
combinations that are semantically incoherent.

\subsection{Cross-Dataset Comparison}

\begin{table}[H]
  \centering
  \caption{Cross-dataset comparison.}
  \label{tab:cross}
  \begin{tabular}{lrr}
    \toprule
    \textbf{Metric} & \textbf{WebQSP} & \textbf{CWQ} \\
    \midrule
    Baseline accuracy & 89.0\% & 53.2\% \\
    Filtering accuracy cost & -1.0pp & -4.0pp \\
    Path reduction & 13.3\% & 10.4\% \\
    Avg paths per question & 2,553 & 1,970 \\
    Avg path length & 2.0 hops & 4.0 hops \\
    \bottomrule
  \end{tabular}
\end{table}

Despite having fewer total paths, CWQ shows higher filtering cost. This is
because CWQ's 4-hop paths form a sparser graph: for a given topic entity and
answer entity, fewer distinct relation sequences connect them. Removing 10\% of
paths is more likely to sever the only correct route.

% ============================================================================
\section{GPU Memory Analysis}
\label{sec:memory}
% ============================================================================

We instrumented the experiment runner to track
\texttt{torch.cuda.max\_memory\_allocated()} per question, measuring peak GPU
memory during trie construction and constrained decoding.

\begin{table}[H]
  \centering
  \caption{Peak GPU memory per question.}
  \label{tab:memory}
  \begin{tabular}{lr}
    \toprule
    \textbf{Method} & \textbf{Avg Peak Memory} \\
    \midrule
    Baseline & 145 MB \\
    Filtered & 151 MB \\
    Label-plan & 148 MB \\
    \bottomrule
  \end{tabular}
\end{table}

The trie construction adds approximately 145--153 MB to the model's base memory
footprint (\(\sim\)16--18 GB for the 8B model in FP16), representing a 0.8--0.9\%
increase---negligible in absolute terms.

Peak memory correlates with the number of \emph{unique tokens} in the path trie,
not with the number of paths directly. Since MarisaTrie compresses shared
prefixes, questions with more diverse entity sets consume more memory regardless
of path count. This explains why reducing path counts by 10--13\% does not
materially reduce memory: the trie's footprint is dominated by entity and
relation token vocabulary.

% ============================================================================
\section{Beam Search vs Greedy Decoding}
\label{sec:beam-vs-greedy}
% ============================================================================

The correction from greedy to beam search (\(k=5\)) is the single largest
accuracy improvement across all methods:

\begin{table}[H]
  \centering
  \caption{Beam vs greedy decoding.}
  \label{tab:beam-vs-greedy}
  \begin{tabular}{lrrr}
    \toprule
    \textbf{Method} & \textbf{Beam (k=5)} & \textbf{Greedy} & \(\Delta\) \\
    \midrule
    Baseline WebQSP & 89.0\% & 80.6\% & +8.4pp \\
    Filtered WebQSP & 88.0\% & 78.9\% & +9.1pp \\
    \bottomrule
  \end{tabular}
\end{table}

Beam search provides a consistent +8--9pp lift regardless of pruning method.
The primary benefit is not finding better paths (the trie restricts paths
identically) but evaluating more diverse completions when the model is uncertain
about intermediate entities. With greedy decoding, a single wrong entity
commitment at hop 1 derails the entire generation.

% ============================================================================
\section{Learned Pruning: A Future Direction}
\label{sec:learned-pruning}
% ============================================================================

The consistent failure of static and heuristic pruning methods motivates a
learned approach: train a fast bi-encoder model to score path relevance and
retain only the top-\(K\) paths before trie construction.

\subsection{Method}

We fine-tune a SentenceTransformer (all-MiniLM-L6-v2, 22M parameters) on
(question, path\_string) pairs. Positive examples are paths whose terminal
entity matches a ground-truth answer; negatives are randomly sampled non-answer
paths. Training uses cosine similarity loss.

The reranker scores all DFS paths for a question in a single forward pass
(embedding the question once, all paths batched). Top-\(K\) paths are retained
for trie construction.

\subsection{Results}

\textbf{Training: 200 questions, evaluation on 81 held-out questions:}

\begin{table}[H]
  \centering
  \caption{Reranker recall@K with 200-sample model.}
  \label{tab:reranker}
  \begin{tabular}{lrrr}
    \toprule
    \textbf{Metric} & \textbf{Reranker} & \textbf{Random} & \textbf{Improvement} \\
    \midrule
    Recall@10 & 45.7\% & 0.5\% & +45.2pp \\
    Recall@50 & 76.2\% & 2.8\% & +73.4pp \\
    Recall@100 & 87.6\% & 5.2\% & +82.4pp \\
    Recall@500 & 98.6\% & 20.7\% & +77.9pp \\
    \bottomrule
  \end{tabular}
\end{table}

\textbf{Zero-shot (no fine-tuning):}

\begin{table}[H]
  \centering
  \caption{Zero-shot reranker recall@K.}
  \label{tab:zeroshot}
  \begin{tabular}{lrrr}
    \toprule
    \textbf{Metric} & \textbf{Cosine Sim} & \textbf{Random} & \textbf{Improvement} \\
    \midrule
    Recall@10 & 27.8\% & 0.3\% & +27.5pp \\
    Recall@100 & 65.0\% & 3.3\% & +61.7pp \\
    Recall@500 & 98.6\% & 17.2\% & +81.4pp \\
    \bottomrule
  \end{tabular}
\end{table}

\subsection{Analysis}

The bi-encoder reranker is effective even zero-shot: simple cosine similarity
between question and path embeddings significantly outperforms random ranking.
Path strings contain entity names and relation names that naturally overlap with
question terms.

The 200-sample model achieves recall@100 = 87.6\%, meaning that with \(K=100\),
we can reduce from \(\sim\)3,400 paths to 100 (97\% reduction) while keeping the
gold path 87.6\% of the time. Combined with GCR's constrained decoding, this
implies an expected end-to-end accuracy of approximately:

\[
  \text{recall@K} \times \text{baseline accuracy} =
  87.6\% \times 53.2\% \approx 46.6\%
\]

This is 6.6pp below the full baseline (53.2\%), but the 97\% path reduction may
improve generation quality by removing distracting alternatives. The key
tradeoff is between path reduction and gold-path retention.

\subsection{Limitations}

\begin{enumerate}
  \item \textbf{No end-to-end validation:} Recall@K measures path ranking only.
        Actual end-to-end accuracy requires running constrained decoding with
        the pruned trie.
  \item \textbf{Training data noise:} Path relevance labels (terminal entity
        matches answer) are noisy. Multiple paths may lead to the same answer
        entity through different relation sequences.
  \item \textbf{Model capacity:} all-MiniLM-L6-v2 (22M parameters) may be too
        small for CWQ's full question diversity. A larger sentence transformer
        may improve recall.
\end{enumerate}

% ============================================================================
\section{Summary of Findings}
\label{sec:summary}
% ============================================================================

\subsection{Main Results}

\begin{table}[H]
  \centering
  \caption{Summary of all methods across datasets.}
  \label{tab:summary}
  \begin{tabular}{lrr}
    \toprule
    \textbf{Method} & \textbf{WebQSP} & \textbf{CWQ} \\
    \midrule
    Baseline (all paths) & \textbf{89.0\%} & \textbf{53.2\%} \\
    Filtered (TypeOracle) & 88.0\% (-1.0pp) & \(\sim\)49\% (-4pp est.) \\
    v2 (step-wise) & 54.9\% (-36.7pp) & 32.2\% (-21.0pp) \\
    Label-reason (10q) & --- & 33.3\% (-19.9pp) \\
    Post-hoc validate & 23.8\% catch rate & --- \\
    Adaptive-500 & 30.7\% & --- \\
    \bottomrule
  \end{tabular}
\end{table}

\subsection{Key Findings}

\begin{enumerate}
  \item \textbf{Static pre-compilation wins:} Baseline 89.0\% WebQSP, 53.2\%
        CWQ---no dynamic method exceeds this.
  \item \textbf{TypeOracle filtering is marginal:} -1pp WebQSP, -4pp CWQ for
        10--13\% path reduction.
  \item \textbf{Step-wise decoding is broken:} -21pp to -37pp; tokenization
        misalignment is a fundamental architectural issue.
  \item \textbf{Ontology reasoning is immature:} 33.3\% on 10 CWQ questions;
        7-category graph too coarse for precision.
  \item \textbf{Beam search is critical:} +8--9pp over greedy; all prior
        results understated.
  \item \textbf{Memory overhead is negligible:} \(\sim\)150 MB for trie
        construction vs 16--18 GB model footprint.
  \item \textbf{Learned pruning is promising:} 87.6\% recall@100 with 97\%
        path reduction; a 22M parameter bi-encoder can dramatically reduce the
        path space while retaining most gold paths.
\end{enumerate}

\subsection{Core Thesis Argument}

Dynamic constraint adaptation for KG-constrained LLM decoding is feasible but
cannot match static pre-compilation accuracy without semantically informed
pruning. The primary value of constrained decoding is hallucination
elimination---the accuracy ceiling is set by the quality of path enumeration,
not by the constraint mechanism itself.

The most promising direction for improvement is learned path reranking: a small
bi-encoder (22M parameters) can reduce the path space by 97\% while retaining
87.6\% of gold paths. This combines the hallucination guarantee of constrained
decoding with the efficiency of query-aware path selection.

% ============================================================================
\section{Legacy Results (Greedy Decoding)}
\label{sec:legacy}
% ============================================================================

These results use greedy decoding (effectively num\_beams=1 due to the
generation config bug). They are included for reference.

\subsection{WebQSP Full Dataset (N=1,627)}

\begin{table}[H]
  \centering
  \caption{WebQSP full dataset results (greedy).}
  \label{tab:legacy-webqsp}
  \begin{tabular}{lrr}
    \toprule
    \textbf{Method} & \textbf{Hits@1} & \textbf{Path Reduction} \\
    \midrule
    Baseline & 80.6\% & --- \\
    Filtered & 78.9\% & 13.3\% \\
    Validate (post-hoc) & 80.6\% & 0\% \\
    Adaptive-500 & 30.7\% & 80\% \\
    Adaptive-100 & 12.8\% & 96\% \\
    Adaptive-30 & 7.9\% & 99\% \\
    \bottomrule
  \end{tabular}
\end{table}

\subsection{CWQ (N=100, Greedy)}

\begin{table}[H]
  \centering
  \caption{CWQ 100q results (greedy).}
  \label{tab:legacy-cwq}
  \begin{tabular}{lrr}
    \toprule
    \textbf{Method} & \textbf{Hits@1} & \(\Delta\) vs Baseline \\
    \midrule
    Baseline & 50.0\% & --- \\
    Filtered & 48.0\% & -2.0pp \\
    v2 & 31.0\% & -19.0pp \\
    Label-plan & 39.0\% & -11.0pp \\
    \bottomrule
  \end{tabular}
\end{table}

% ============================================================================
\section{Limitations}
\label{sec:limitations}
% ============================================================================

\begin{enumerate}
  \item \textbf{WebQSP beam run:} 580 of 1,628 samples completed under beam
        search. Full-dataset confirmation would increase statistical confidence.
  \item \textbf{CWQ filtered at 500q scale:} Only 100q filtered results are
        available experimentally. The -4pp delta observed at 100q may not fully
        generalize, though the direction is consistent.
  \item \textbf{Label-reason evaluation:} Only 10 CWQ questions evaluated. The
        33.3\% result is indicative but not statistically robust.
  \item \textbf{No small-model comparison:} The hypothesis that a smaller model
        (e.g., Qwen 3B) would benefit more from TypeOracle pruning could not be
        tested due to GPU constraints.
  \item \textbf{Adaptive budget study:} Current adaptive methods use DFS-order
        path truncation, introducing strong selection bias. Results are not
        representative of what relevance-based path selection would achieve.
  \item \textbf{Learned pruning:} No end-to-end validation. Recall@K measures
        path ranking quality but does not account for the interaction between
        path pruning and constrained decoding behavior.
\end{enumerate}

% ============================================================================
\section{Related Work}
\label{sec:related}
% ============================================================================

\textbf{Graph-Constrained Decoding.} GCR~\citep{luo2025-graph-constrained-reasoning}
pre-compiles KG paths into a trie for prefix-constrained generation.
DoG~\citep{li2024-decoding-on-graphs} extends this to dynamic graph contexts.
Our work replicates and extends GCR with six dynamic constraint strategies.

\textbf{Type-Based Filtering.} TypeOracle~\citep{ravichander2022typeoracle}
uses Freebase type annotations to prune implausible paths.
TypeGate~\citep{fernando2023typegate} applies type constraints during decoding.
Our Filtered and Validate methods implement these ideas within the trie
framework.

\textbf{Step-Wise Reasoning.} StructGPT~\citep{jiang2023structgpt} and
Interleaved Retrieval~\citep{feng2024interleaved} decompose multi-hop reasoning
into sequential KG queries. Our v2 method follows this paradigm but reveals
fundamental tokenization alignment issues.

\textbf{Ontology Reasoning.} ORT~\citep{ort2025} plans at the label level before
entity enumeration. Our label-reason implementation confirms the approach's
feasibility but shows that auto-mined ontologies lack the precision of
hand-curated knowledge.

\textbf{Learned Path Ranking.} REL-RAG~\citep{tang2024relrag} uses a retriever
to select relevant paths before generation. REAP~\citep{zhang2025reap} jointly
learns path selection and answer generation. Our bi-encoder reranker is closest
to REL-RAG's retrieval stage.

% ============================================================================
\bibliographystyle{plainnat}
\bibliography{refs}
% ============================================================================

\end{document}
